\begin{abstract}
\textit{Cross-view object correspondence} is a fundamental task for multi-camera intelligent vision systems, requiring a method to identify the same physical object across cameras despite drastic differences in scale, lighting, and visibility and in the absence of geometric calibration. \textit{EgoExo-4D}, the primary benchmark for this task, has attracted a growing body of prior work, and existing methods achieve respectable performance. However, every published method relies on an architecture trained directly on benchmark data, making it inapplicable to novel camera configurations and opaque to inspection. We ask whether natural language can serve as a bridge between machine learning architectures, compressing a scene into a human-readable description that transfers object identity from one view to another. This framing is motivated by how humans collaborate through language and by the interpretability such a bridge affords, and it subsumes a second question of whether general-purpose foundation models can collaborate, composed at inference time, on a task none was trained to solve.

We present a fully training-free pipeline in which natural language serves as the bridge between viewpoints. A vision-language model describes the target object using view- and time-independent terms; an open-set detector grounds that description to select anchor frames in the destination view; \textit{Segment Anything Model~3 Agent} segments and propagates the mask across frames. No component is fine-tuned, and the bridge between views is a human-readable phrase that renders every intermediate decision inspectable without accessing model weights. We benchmark our approach against published methods, reaching within a few percentage points of the 2025 challenge winners, and provide a detailed ablation study examining where foundation models succeed and where they remain limited on a task none were designed for.

The thesis develops these contributions across six chapters. A theoretical chapter first establishes the mechanistic foundations of each pipeline component, covering visual representation learning with self-supervised transformers, image and video segmentation with the Segment Anything family, language-grounded open-set detection, large language and vision-language models, and temporal mask propagation. We then survey related work and formally define the correspondence task, present the proposed pipeline with motivation for each design choice, and close with experimental results, ablation analyses, and a discussion of limitations and future directions. Together, these chapters demonstrate that a language-mediated composition of general-purpose foundation models constitutes a viable, interpretable, and training-free alternative for cross-view object correspondence, a task previously dominated by specialised learned architectures.
\end{abstract}
